% ============================================================================
% ESQUEMA PROPOSTO (banca, sugestão — não mergeado): o algoritmo DRI-SL como
% fluxo em duas fases — (i) densidade semântica e (ii) variedade lexical
% intragrupo — seguido da avaliação. Ilustra a Seção sec:metodo-drisl.
%
% Uso no Cap. 3 (dentro de um ambiente figure):
%   % ============================================================================
% ESQUEMA PROPOSTO (banca, sugestão — não mergeado): o algoritmo DRI-SL como
% fluxo em duas fases — (i) densidade semântica e (ii) variedade lexical
% intragrupo — seguido da avaliação. Ilustra a Seção sec:metodo-drisl.
%
% Uso no Cap. 3 (dentro de um ambiente figure):
%   % ============================================================================
% ESQUEMA PROPOSTO (banca, sugestão — não mergeado): o algoritmo DRI-SL como
% fluxo em duas fases — (i) densidade semântica e (ii) variedade lexical
% intragrupo — seguido da avaliação. Ilustra a Seção sec:metodo-drisl.
%
% Uso no Cap. 3 (dentro de um ambiente figure):
%   % ============================================================================
% ESQUEMA PROPOSTO (banca, sugestão — não mergeado): o algoritmo DRI-SL como
% fluxo em duas fases — (i) densidade semântica e (ii) variedade lexical
% intragrupo — seguido da avaliação. Ilustra a Seção sec:metodo-drisl.
%
% Uso no Cap. 3 (dentro de um ambiente figure):
%   \input{3-metodo/esquemas-propostos/esq-drisl}
%   A citação do codificador de sentenças (\citep{Reimers2019SBERT}) vai na
%   LEGENDA da figure, não pintada no desenho — assim o check-bib a enxerga.
% Uso standalone (pré-visualização):
%   \documentclass[tikz,border=6pt]{standalone}
%   \usetikzlibrary{arrows.meta, positioning}
%   \begin{document}\input{esq-drisl.tex}\end{document}
%   (no modo standalone, troque as \ref{...} por texto fixo ou defina-as)
% Requer apenas arrows.meta e positioning (já em packages.tex). Legível em
% P&B: etapa de medição/entrada = cinza claro; processamento = branco;
% produto = cinza médio; remissão = texto simples.
% FREEZE respeitado: nenhum número de resultado; só os símbolos da própria
% seção (U0, L0, I, Nc, Gi, ni) e nenhum gráfico decorativo que pareça dado.
% Terminologia em camadas: a figura não afirma o resultado da comparação
% (isso é do Capítulo 4); só diz CONTRA O QUE o DRI-SL é avaliado.
% ============================================================================
\begin{tikzpicture}[
    x=1cm, y=1cm,
    caixa/.style={draw, rounded corners=2pt, align=center, text width=34mm,
                  minimum height=11mm, inner sep=3pt, font=\footnotesize},
    entrada/.style={caixa, fill=gray!8},
    produto/.style={caixa, fill=gray!14},
    seta/.style={-{Stealth[length=2.4mm]}, thick},
    rot/.style={font=\scriptsize, align=center, inner sep=2pt}
]
    % ---------------- fase (i): densidade semântica ------------------------
    \node[rot, anchor=west] (tit1) at (-1.9,1.1)
        {\textbf{Fase (i) --- densidade semântica}};

    \node[entrada] (u0) at (0,0)
        {conjunto não rotulado $U_0$ (textos curtos)};
    \node[caixa, text width=36mm] (enc) at (4.6,0)
        {codificador de\\sentenças: projeção\\no espaço semântico};
    \node[caixa] (km) at (9.2,0)
        {$k$-médias em $N_c$ grupos semânticos};
    \draw[seta] (u0) -- (enc);
    \draw[seta] (enc) -- (km);

    \node[produto] (aloc) at (9.2,-2.3)
        {cotas proporcionais ao tamanho de cada grupo:
         $\sum_{i=1}^{N_c} n_i = I$};
    \draw[seta] (km) -- (aloc);

    % ---------------- fase (ii): variedade lexical intragrupo --------------
    % o título da fase É o rótulo da seta em L (um elemento só: título de
    % banda separado era atravessado pela descida da seta — achado visual)
    \node[caixa] (cand) at (0,-5.4)
        {candidatos do grupo $G_i$ (ainda não selecionados)};
    \node[caixa, text width=40mm] (sel) at (4.6,-5.4)
        {seleciona a instância de maior escore de novidade, ponderado pelo
         perfil TF--IDF do grupo};
    \node[produto, text width=44mm] (l0) at (9.6,-5.4)
        {conjunto inicial $L_0$ (\mbox{$|L_0|=I$}):\\
         diversidade semântica entre\\
         grupos $+$ variedade lexical\\
         dentro de cada grupo};

    % a linha em L conduz da alocação para o laço de seleção; o rótulo na
    % horizontal é o título da fase (desenhada aqui, depois dos nós)
    \draw[seta] (aloc.south) -- ++(0,-0.55) -| (cand.north);
    % rótulo-título posicionado à mão sobre o trecho horizontal, com as
    % bordas afastadas das duas descidas verticais (x=0 e x=9,2)
    \node[rot, anchor=south] at (4.2,-3.42)
        {\textbf{Fase (ii) --- variedade lexical intragrupo}:\\
         para cada grupo $G_i$, com cota $n_i$};

    \draw[seta] (cand) -- (sel);
    \draw[seta] (sel) -- (l0);
    \draw[seta] (sel.south) to[out=250, in=290]
        node[rot, below=2pt] {atualiza os termos cobertos; repete até $n_i$ instâncias}
        (cand.south);

    % ---------------- avaliação (sem afirmar resultado) --------------------
    \node[rot, anchor=west, text width=124mm, align=left] (tit3) at (-1.9,-8.6)
        {\textbf{Avaliação} (resultados na Seção~\ref{sec:res-l0-drisl}): o
         $L_0$ do DRI-SL é comparado ao da amostragem aleatória e ao envelope
         de desempenho do algoritmo genético, que atua como limite empírico
         de referência.};
\end{tikzpicture}

%   A citação do codificador de sentenças (\citep{Reimers2019SBERT}) vai na
%   LEGENDA da figure, não pintada no desenho — assim o check-bib a enxerga.
% Uso standalone (pré-visualização):
%   \documentclass[tikz,border=6pt]{standalone}
%   \usetikzlibrary{arrows.meta, positioning}
%   \begin{document}% ============================================================================
% ESQUEMA PROPOSTO (banca, sugestão — não mergeado): o algoritmo DRI-SL como
% fluxo em duas fases — (i) densidade semântica e (ii) variedade lexical
% intragrupo — seguido da avaliação. Ilustra a Seção sec:metodo-drisl.
%
% Uso no Cap. 3 (dentro de um ambiente figure):
%   \input{3-metodo/esquemas-propostos/esq-drisl}
%   A citação do codificador de sentenças (\citep{Reimers2019SBERT}) vai na
%   LEGENDA da figure, não pintada no desenho — assim o check-bib a enxerga.
% Uso standalone (pré-visualização):
%   \documentclass[tikz,border=6pt]{standalone}
%   \usetikzlibrary{arrows.meta, positioning}
%   \begin{document}\input{esq-drisl.tex}\end{document}
%   (no modo standalone, troque as \ref{...} por texto fixo ou defina-as)
% Requer apenas arrows.meta e positioning (já em packages.tex). Legível em
% P&B: etapa de medição/entrada = cinza claro; processamento = branco;
% produto = cinza médio; remissão = texto simples.
% FREEZE respeitado: nenhum número de resultado; só os símbolos da própria
% seção (U0, L0, I, Nc, Gi, ni) e nenhum gráfico decorativo que pareça dado.
% Terminologia em camadas: a figura não afirma o resultado da comparação
% (isso é do Capítulo 4); só diz CONTRA O QUE o DRI-SL é avaliado.
% ============================================================================
\begin{tikzpicture}[
    x=1cm, y=1cm,
    caixa/.style={draw, rounded corners=2pt, align=center, text width=34mm,
                  minimum height=11mm, inner sep=3pt, font=\footnotesize},
    entrada/.style={caixa, fill=gray!8},
    produto/.style={caixa, fill=gray!14},
    seta/.style={-{Stealth[length=2.4mm]}, thick},
    rot/.style={font=\scriptsize, align=center, inner sep=2pt}
]
    % ---------------- fase (i): densidade semântica ------------------------
    \node[rot, anchor=west] (tit1) at (-1.9,1.1)
        {\textbf{Fase (i) --- densidade semântica}};

    \node[entrada] (u0) at (0,0)
        {conjunto não rotulado $U_0$ (textos curtos)};
    \node[caixa, text width=36mm] (enc) at (4.6,0)
        {codificador de\\sentenças: projeção\\no espaço semântico};
    \node[caixa] (km) at (9.2,0)
        {$k$-médias em $N_c$ grupos semânticos};
    \draw[seta] (u0) -- (enc);
    \draw[seta] (enc) -- (km);

    \node[produto] (aloc) at (9.2,-2.3)
        {cotas proporcionais ao tamanho de cada grupo:
         $\sum_{i=1}^{N_c} n_i = I$};
    \draw[seta] (km) -- (aloc);

    % ---------------- fase (ii): variedade lexical intragrupo --------------
    % o título da fase É o rótulo da seta em L (um elemento só: título de
    % banda separado era atravessado pela descida da seta — achado visual)
    \node[caixa] (cand) at (0,-5.4)
        {candidatos do grupo $G_i$ (ainda não selecionados)};
    \node[caixa, text width=40mm] (sel) at (4.6,-5.4)
        {seleciona a instância de maior escore de novidade, ponderado pelo
         perfil TF--IDF do grupo};
    \node[produto, text width=44mm] (l0) at (9.6,-5.4)
        {conjunto inicial $L_0$ (\mbox{$|L_0|=I$}):\\
         diversidade semântica entre\\
         grupos $+$ variedade lexical\\
         dentro de cada grupo};

    % a linha em L conduz da alocação para o laço de seleção; o rótulo na
    % horizontal é o título da fase (desenhada aqui, depois dos nós)
    \draw[seta] (aloc.south) -- ++(0,-0.55) -| (cand.north);
    % rótulo-título posicionado à mão sobre o trecho horizontal, com as
    % bordas afastadas das duas descidas verticais (x=0 e x=9,2)
    \node[rot, anchor=south] at (4.2,-3.42)
        {\textbf{Fase (ii) --- variedade lexical intragrupo}:\\
         para cada grupo $G_i$, com cota $n_i$};

    \draw[seta] (cand) -- (sel);
    \draw[seta] (sel) -- (l0);
    \draw[seta] (sel.south) to[out=250, in=290]
        node[rot, below=2pt] {atualiza os termos cobertos; repete até $n_i$ instâncias}
        (cand.south);

    % ---------------- avaliação (sem afirmar resultado) --------------------
    \node[rot, anchor=west, text width=124mm, align=left] (tit3) at (-1.9,-8.6)
        {\textbf{Avaliação} (resultados na Seção~\ref{sec:res-l0-drisl}): o
         $L_0$ do DRI-SL é comparado ao da amostragem aleatória e ao envelope
         de desempenho do algoritmo genético, que atua como limite empírico
         de referência.};
\end{tikzpicture}
\end{document}
%   (no modo standalone, troque as \ref{...} por texto fixo ou defina-as)
% Requer apenas arrows.meta e positioning (já em packages.tex). Legível em
% P&B: etapa de medição/entrada = cinza claro; processamento = branco;
% produto = cinza médio; remissão = texto simples.
% FREEZE respeitado: nenhum número de resultado; só os símbolos da própria
% seção (U0, L0, I, Nc, Gi, ni) e nenhum gráfico decorativo que pareça dado.
% Terminologia em camadas: a figura não afirma o resultado da comparação
% (isso é do Capítulo 4); só diz CONTRA O QUE o DRI-SL é avaliado.
% ============================================================================
\begin{tikzpicture}[
    x=1cm, y=1cm,
    caixa/.style={draw, rounded corners=2pt, align=center, text width=34mm,
                  minimum height=11mm, inner sep=3pt, font=\footnotesize},
    entrada/.style={caixa, fill=gray!8},
    produto/.style={caixa, fill=gray!14},
    seta/.style={-{Stealth[length=2.4mm]}, thick},
    rot/.style={font=\scriptsize, align=center, inner sep=2pt}
]
    % ---------------- fase (i): densidade semântica ------------------------
    \node[rot, anchor=west] (tit1) at (-1.9,1.1)
        {\textbf{Fase (i) --- densidade semântica}};

    \node[entrada] (u0) at (0,0)
        {conjunto não rotulado $U_0$ (textos curtos)};
    \node[caixa, text width=36mm] (enc) at (4.6,0)
        {codificador de\\sentenças: projeção\\no espaço semântico};
    \node[caixa] (km) at (9.2,0)
        {$k$-médias em $N_c$ grupos semânticos};
    \draw[seta] (u0) -- (enc);
    \draw[seta] (enc) -- (km);

    \node[produto] (aloc) at (9.2,-2.3)
        {cotas proporcionais ao tamanho de cada grupo:
         $\sum_{i=1}^{N_c} n_i = I$};
    \draw[seta] (km) -- (aloc);

    % ---------------- fase (ii): variedade lexical intragrupo --------------
    % o título da fase É o rótulo da seta em L (um elemento só: título de
    % banda separado era atravessado pela descida da seta — achado visual)
    \node[caixa] (cand) at (0,-5.4)
        {candidatos do grupo $G_i$ (ainda não selecionados)};
    \node[caixa, text width=40mm] (sel) at (4.6,-5.4)
        {seleciona a instância de maior escore de novidade, ponderado pelo
         perfil TF--IDF do grupo};
    \node[produto, text width=44mm] (l0) at (9.6,-5.4)
        {conjunto inicial $L_0$ (\mbox{$|L_0|=I$}):\\
         diversidade semântica entre\\
         grupos $+$ variedade lexical\\
         dentro de cada grupo};

    % a linha em L conduz da alocação para o laço de seleção; o rótulo na
    % horizontal é o título da fase (desenhada aqui, depois dos nós)
    \draw[seta] (aloc.south) -- ++(0,-0.55) -| (cand.north);
    % rótulo-título posicionado à mão sobre o trecho horizontal, com as
    % bordas afastadas das duas descidas verticais (x=0 e x=9,2)
    \node[rot, anchor=south] at (4.2,-3.42)
        {\textbf{Fase (ii) --- variedade lexical intragrupo}:\\
         para cada grupo $G_i$, com cota $n_i$};

    \draw[seta] (cand) -- (sel);
    \draw[seta] (sel) -- (l0);
    \draw[seta] (sel.south) to[out=250, in=290]
        node[rot, below=2pt] {atualiza os termos cobertos; repete até $n_i$ instâncias}
        (cand.south);

    % ---------------- avaliação (sem afirmar resultado) --------------------
    \node[rot, anchor=west, text width=124mm, align=left] (tit3) at (-1.9,-8.6)
        {\textbf{Avaliação} (resultados na Seção~\ref{sec:res-l0-drisl}): o
         $L_0$ do DRI-SL é comparado ao da amostragem aleatória e ao envelope
         de desempenho do algoritmo genético, que atua como limite empírico
         de referência.};
\end{tikzpicture}

%   A citação do codificador de sentenças (\citep{Reimers2019SBERT}) vai na
%   LEGENDA da figure, não pintada no desenho — assim o check-bib a enxerga.
% Uso standalone (pré-visualização):
%   \documentclass[tikz,border=6pt]{standalone}
%   \usetikzlibrary{arrows.meta, positioning}
%   \begin{document}% ============================================================================
% ESQUEMA PROPOSTO (banca, sugestão — não mergeado): o algoritmo DRI-SL como
% fluxo em duas fases — (i) densidade semântica e (ii) variedade lexical
% intragrupo — seguido da avaliação. Ilustra a Seção sec:metodo-drisl.
%
% Uso no Cap. 3 (dentro de um ambiente figure):
%   % ============================================================================
% ESQUEMA PROPOSTO (banca, sugestão — não mergeado): o algoritmo DRI-SL como
% fluxo em duas fases — (i) densidade semântica e (ii) variedade lexical
% intragrupo — seguido da avaliação. Ilustra a Seção sec:metodo-drisl.
%
% Uso no Cap. 3 (dentro de um ambiente figure):
%   \input{3-metodo/esquemas-propostos/esq-drisl}
%   A citação do codificador de sentenças (\citep{Reimers2019SBERT}) vai na
%   LEGENDA da figure, não pintada no desenho — assim o check-bib a enxerga.
% Uso standalone (pré-visualização):
%   \documentclass[tikz,border=6pt]{standalone}
%   \usetikzlibrary{arrows.meta, positioning}
%   \begin{document}\input{esq-drisl.tex}\end{document}
%   (no modo standalone, troque as \ref{...} por texto fixo ou defina-as)
% Requer apenas arrows.meta e positioning (já em packages.tex). Legível em
% P&B: etapa de medição/entrada = cinza claro; processamento = branco;
% produto = cinza médio; remissão = texto simples.
% FREEZE respeitado: nenhum número de resultado; só os símbolos da própria
% seção (U0, L0, I, Nc, Gi, ni) e nenhum gráfico decorativo que pareça dado.
% Terminologia em camadas: a figura não afirma o resultado da comparação
% (isso é do Capítulo 4); só diz CONTRA O QUE o DRI-SL é avaliado.
% ============================================================================
\begin{tikzpicture}[
    x=1cm, y=1cm,
    caixa/.style={draw, rounded corners=2pt, align=center, text width=34mm,
                  minimum height=11mm, inner sep=3pt, font=\footnotesize},
    entrada/.style={caixa, fill=gray!8},
    produto/.style={caixa, fill=gray!14},
    seta/.style={-{Stealth[length=2.4mm]}, thick},
    rot/.style={font=\scriptsize, align=center, inner sep=2pt}
]
    % ---------------- fase (i): densidade semântica ------------------------
    \node[rot, anchor=west] (tit1) at (-1.9,1.1)
        {\textbf{Fase (i) --- densidade semântica}};

    \node[entrada] (u0) at (0,0)
        {conjunto não rotulado $U_0$ (textos curtos)};
    \node[caixa, text width=36mm] (enc) at (4.6,0)
        {codificador de\\sentenças: projeção\\no espaço semântico};
    \node[caixa] (km) at (9.2,0)
        {$k$-médias em $N_c$ grupos semânticos};
    \draw[seta] (u0) -- (enc);
    \draw[seta] (enc) -- (km);

    \node[produto] (aloc) at (9.2,-2.3)
        {cotas proporcionais ao tamanho de cada grupo:
         $\sum_{i=1}^{N_c} n_i = I$};
    \draw[seta] (km) -- (aloc);

    % ---------------- fase (ii): variedade lexical intragrupo --------------
    % o título da fase É o rótulo da seta em L (um elemento só: título de
    % banda separado era atravessado pela descida da seta — achado visual)
    \node[caixa] (cand) at (0,-5.4)
        {candidatos do grupo $G_i$ (ainda não selecionados)};
    \node[caixa, text width=40mm] (sel) at (4.6,-5.4)
        {seleciona a instância de maior escore de novidade, ponderado pelo
         perfil TF--IDF do grupo};
    \node[produto, text width=44mm] (l0) at (9.6,-5.4)
        {conjunto inicial $L_0$ (\mbox{$|L_0|=I$}):\\
         diversidade semântica entre\\
         grupos $+$ variedade lexical\\
         dentro de cada grupo};

    % a linha em L conduz da alocação para o laço de seleção; o rótulo na
    % horizontal é o título da fase (desenhada aqui, depois dos nós)
    \draw[seta] (aloc.south) -- ++(0,-0.55) -| (cand.north);
    % rótulo-título posicionado à mão sobre o trecho horizontal, com as
    % bordas afastadas das duas descidas verticais (x=0 e x=9,2)
    \node[rot, anchor=south] at (4.2,-3.42)
        {\textbf{Fase (ii) --- variedade lexical intragrupo}:\\
         para cada grupo $G_i$, com cota $n_i$};

    \draw[seta] (cand) -- (sel);
    \draw[seta] (sel) -- (l0);
    \draw[seta] (sel.south) to[out=250, in=290]
        node[rot, below=2pt] {atualiza os termos cobertos; repete até $n_i$ instâncias}
        (cand.south);

    % ---------------- avaliação (sem afirmar resultado) --------------------
    \node[rot, anchor=west, text width=124mm, align=left] (tit3) at (-1.9,-8.6)
        {\textbf{Avaliação} (resultados na Seção~\ref{sec:res-l0-drisl}): o
         $L_0$ do DRI-SL é comparado ao da amostragem aleatória e ao envelope
         de desempenho do algoritmo genético, que atua como limite empírico
         de referência.};
\end{tikzpicture}

%   A citação do codificador de sentenças (\citep{Reimers2019SBERT}) vai na
%   LEGENDA da figure, não pintada no desenho — assim o check-bib a enxerga.
% Uso standalone (pré-visualização):
%   \documentclass[tikz,border=6pt]{standalone}
%   \usetikzlibrary{arrows.meta, positioning}
%   \begin{document}% ============================================================================
% ESQUEMA PROPOSTO (banca, sugestão — não mergeado): o algoritmo DRI-SL como
% fluxo em duas fases — (i) densidade semântica e (ii) variedade lexical
% intragrupo — seguido da avaliação. Ilustra a Seção sec:metodo-drisl.
%
% Uso no Cap. 3 (dentro de um ambiente figure):
%   \input{3-metodo/esquemas-propostos/esq-drisl}
%   A citação do codificador de sentenças (\citep{Reimers2019SBERT}) vai na
%   LEGENDA da figure, não pintada no desenho — assim o check-bib a enxerga.
% Uso standalone (pré-visualização):
%   \documentclass[tikz,border=6pt]{standalone}
%   \usetikzlibrary{arrows.meta, positioning}
%   \begin{document}\input{esq-drisl.tex}\end{document}
%   (no modo standalone, troque as \ref{...} por texto fixo ou defina-as)
% Requer apenas arrows.meta e positioning (já em packages.tex). Legível em
% P&B: etapa de medição/entrada = cinza claro; processamento = branco;
% produto = cinza médio; remissão = texto simples.
% FREEZE respeitado: nenhum número de resultado; só os símbolos da própria
% seção (U0, L0, I, Nc, Gi, ni) e nenhum gráfico decorativo que pareça dado.
% Terminologia em camadas: a figura não afirma o resultado da comparação
% (isso é do Capítulo 4); só diz CONTRA O QUE o DRI-SL é avaliado.
% ============================================================================
\begin{tikzpicture}[
    x=1cm, y=1cm,
    caixa/.style={draw, rounded corners=2pt, align=center, text width=34mm,
                  minimum height=11mm, inner sep=3pt, font=\footnotesize},
    entrada/.style={caixa, fill=gray!8},
    produto/.style={caixa, fill=gray!14},
    seta/.style={-{Stealth[length=2.4mm]}, thick},
    rot/.style={font=\scriptsize, align=center, inner sep=2pt}
]
    % ---------------- fase (i): densidade semântica ------------------------
    \node[rot, anchor=west] (tit1) at (-1.9,1.1)
        {\textbf{Fase (i) --- densidade semântica}};

    \node[entrada] (u0) at (0,0)
        {conjunto não rotulado $U_0$ (textos curtos)};
    \node[caixa, text width=36mm] (enc) at (4.6,0)
        {codificador de\\sentenças: projeção\\no espaço semântico};
    \node[caixa] (km) at (9.2,0)
        {$k$-médias em $N_c$ grupos semânticos};
    \draw[seta] (u0) -- (enc);
    \draw[seta] (enc) -- (km);

    \node[produto] (aloc) at (9.2,-2.3)
        {cotas proporcionais ao tamanho de cada grupo:
         $\sum_{i=1}^{N_c} n_i = I$};
    \draw[seta] (km) -- (aloc);

    % ---------------- fase (ii): variedade lexical intragrupo --------------
    % o título da fase É o rótulo da seta em L (um elemento só: título de
    % banda separado era atravessado pela descida da seta — achado visual)
    \node[caixa] (cand) at (0,-5.4)
        {candidatos do grupo $G_i$ (ainda não selecionados)};
    \node[caixa, text width=40mm] (sel) at (4.6,-5.4)
        {seleciona a instância de maior escore de novidade, ponderado pelo
         perfil TF--IDF do grupo};
    \node[produto, text width=44mm] (l0) at (9.6,-5.4)
        {conjunto inicial $L_0$ (\mbox{$|L_0|=I$}):\\
         diversidade semântica entre\\
         grupos $+$ variedade lexical\\
         dentro de cada grupo};

    % a linha em L conduz da alocação para o laço de seleção; o rótulo na
    % horizontal é o título da fase (desenhada aqui, depois dos nós)
    \draw[seta] (aloc.south) -- ++(0,-0.55) -| (cand.north);
    % rótulo-título posicionado à mão sobre o trecho horizontal, com as
    % bordas afastadas das duas descidas verticais (x=0 e x=9,2)
    \node[rot, anchor=south] at (4.2,-3.42)
        {\textbf{Fase (ii) --- variedade lexical intragrupo}:\\
         para cada grupo $G_i$, com cota $n_i$};

    \draw[seta] (cand) -- (sel);
    \draw[seta] (sel) -- (l0);
    \draw[seta] (sel.south) to[out=250, in=290]
        node[rot, below=2pt] {atualiza os termos cobertos; repete até $n_i$ instâncias}
        (cand.south);

    % ---------------- avaliação (sem afirmar resultado) --------------------
    \node[rot, anchor=west, text width=124mm, align=left] (tit3) at (-1.9,-8.6)
        {\textbf{Avaliação} (resultados na Seção~\ref{sec:res-l0-drisl}): o
         $L_0$ do DRI-SL é comparado ao da amostragem aleatória e ao envelope
         de desempenho do algoritmo genético, que atua como limite empírico
         de referência.};
\end{tikzpicture}
\end{document}
%   (no modo standalone, troque as \ref{...} por texto fixo ou defina-as)
% Requer apenas arrows.meta e positioning (já em packages.tex). Legível em
% P&B: etapa de medição/entrada = cinza claro; processamento = branco;
% produto = cinza médio; remissão = texto simples.
% FREEZE respeitado: nenhum número de resultado; só os símbolos da própria
% seção (U0, L0, I, Nc, Gi, ni) e nenhum gráfico decorativo que pareça dado.
% Terminologia em camadas: a figura não afirma o resultado da comparação
% (isso é do Capítulo 4); só diz CONTRA O QUE o DRI-SL é avaliado.
% ============================================================================
\begin{tikzpicture}[
    x=1cm, y=1cm,
    caixa/.style={draw, rounded corners=2pt, align=center, text width=34mm,
                  minimum height=11mm, inner sep=3pt, font=\footnotesize},
    entrada/.style={caixa, fill=gray!8},
    produto/.style={caixa, fill=gray!14},
    seta/.style={-{Stealth[length=2.4mm]}, thick},
    rot/.style={font=\scriptsize, align=center, inner sep=2pt}
]
    % ---------------- fase (i): densidade semântica ------------------------
    \node[rot, anchor=west] (tit1) at (-1.9,1.1)
        {\textbf{Fase (i) --- densidade semântica}};

    \node[entrada] (u0) at (0,0)
        {conjunto não rotulado $U_0$ (textos curtos)};
    \node[caixa, text width=36mm] (enc) at (4.6,0)
        {codificador de\\sentenças: projeção\\no espaço semântico};
    \node[caixa] (km) at (9.2,0)
        {$k$-médias em $N_c$ grupos semânticos};
    \draw[seta] (u0) -- (enc);
    \draw[seta] (enc) -- (km);

    \node[produto] (aloc) at (9.2,-2.3)
        {cotas proporcionais ao tamanho de cada grupo:
         $\sum_{i=1}^{N_c} n_i = I$};
    \draw[seta] (km) -- (aloc);

    % ---------------- fase (ii): variedade lexical intragrupo --------------
    % o título da fase É o rótulo da seta em L (um elemento só: título de
    % banda separado era atravessado pela descida da seta — achado visual)
    \node[caixa] (cand) at (0,-5.4)
        {candidatos do grupo $G_i$ (ainda não selecionados)};
    \node[caixa, text width=40mm] (sel) at (4.6,-5.4)
        {seleciona a instância de maior escore de novidade, ponderado pelo
         perfil TF--IDF do grupo};
    \node[produto, text width=44mm] (l0) at (9.6,-5.4)
        {conjunto inicial $L_0$ (\mbox{$|L_0|=I$}):\\
         diversidade semântica entre\\
         grupos $+$ variedade lexical\\
         dentro de cada grupo};

    % a linha em L conduz da alocação para o laço de seleção; o rótulo na
    % horizontal é o título da fase (desenhada aqui, depois dos nós)
    \draw[seta] (aloc.south) -- ++(0,-0.55) -| (cand.north);
    % rótulo-título posicionado à mão sobre o trecho horizontal, com as
    % bordas afastadas das duas descidas verticais (x=0 e x=9,2)
    \node[rot, anchor=south] at (4.2,-3.42)
        {\textbf{Fase (ii) --- variedade lexical intragrupo}:\\
         para cada grupo $G_i$, com cota $n_i$};

    \draw[seta] (cand) -- (sel);
    \draw[seta] (sel) -- (l0);
    \draw[seta] (sel.south) to[out=250, in=290]
        node[rot, below=2pt] {atualiza os termos cobertos; repete até $n_i$ instâncias}
        (cand.south);

    % ---------------- avaliação (sem afirmar resultado) --------------------
    \node[rot, anchor=west, text width=124mm, align=left] (tit3) at (-1.9,-8.6)
        {\textbf{Avaliação} (resultados na Seção~\ref{sec:res-l0-drisl}): o
         $L_0$ do DRI-SL é comparado ao da amostragem aleatória e ao envelope
         de desempenho do algoritmo genético, que atua como limite empírico
         de referência.};
\end{tikzpicture}
\end{document}
%   (no modo standalone, troque as \ref{...} por texto fixo ou defina-as)
% Requer apenas arrows.meta e positioning (já em packages.tex). Legível em
% P&B: etapa de medição/entrada = cinza claro; processamento = branco;
% produto = cinza médio; remissão = texto simples.
% FREEZE respeitado: nenhum número de resultado; só os símbolos da própria
% seção (U0, L0, I, Nc, Gi, ni) e nenhum gráfico decorativo que pareça dado.
% Terminologia em camadas: a figura não afirma o resultado da comparação
% (isso é do Capítulo 4); só diz CONTRA O QUE o DRI-SL é avaliado.
% ============================================================================
\begin{tikzpicture}[
    x=1cm, y=1cm,
    caixa/.style={draw, rounded corners=2pt, align=center, text width=34mm,
                  minimum height=11mm, inner sep=3pt, font=\footnotesize},
    entrada/.style={caixa, fill=gray!8},
    produto/.style={caixa, fill=gray!14},
    seta/.style={-{Stealth[length=2.4mm]}, thick},
    rot/.style={font=\scriptsize, align=center, inner sep=2pt}
]
    % ---------------- fase (i): densidade semântica ------------------------
    \node[rot, anchor=west] (tit1) at (-1.9,1.1)
        {\textbf{Fase (i) --- densidade semântica}};

    \node[entrada] (u0) at (0,0)
        {conjunto não rotulado $U_0$ (textos curtos)};
    \node[caixa, text width=36mm] (enc) at (4.6,0)
        {codificador de\\sentenças: projeção\\no espaço semântico};
    \node[caixa] (km) at (9.2,0)
        {$k$-médias em $N_c$ grupos semânticos};
    \draw[seta] (u0) -- (enc);
    \draw[seta] (enc) -- (km);

    \node[produto] (aloc) at (9.2,-2.3)
        {cotas proporcionais ao tamanho de cada grupo:
         $\sum_{i=1}^{N_c} n_i = I$};
    \draw[seta] (km) -- (aloc);

    % ---------------- fase (ii): variedade lexical intragrupo --------------
    % o título da fase É o rótulo da seta em L (um elemento só: título de
    % banda separado era atravessado pela descida da seta — achado visual)
    \node[caixa] (cand) at (0,-5.4)
        {candidatos do grupo $G_i$ (ainda não selecionados)};
    \node[caixa, text width=40mm] (sel) at (4.6,-5.4)
        {seleciona a instância de maior escore de novidade, ponderado pelo
         perfil TF--IDF do grupo};
    \node[produto, text width=44mm] (l0) at (9.6,-5.4)
        {conjunto inicial $L_0$ (\mbox{$|L_0|=I$}):\\
         diversidade semântica entre\\
         grupos $+$ variedade lexical\\
         dentro de cada grupo};

    % a linha em L conduz da alocação para o laço de seleção; o rótulo na
    % horizontal é o título da fase (desenhada aqui, depois dos nós)
    \draw[seta] (aloc.south) -- ++(0,-0.55) -| (cand.north);
    % rótulo-título posicionado à mão sobre o trecho horizontal, com as
    % bordas afastadas das duas descidas verticais (x=0 e x=9,2)
    \node[rot, anchor=south] at (4.2,-3.42)
        {\textbf{Fase (ii) --- variedade lexical intragrupo}:\\
         para cada grupo $G_i$, com cota $n_i$};

    \draw[seta] (cand) -- (sel);
    \draw[seta] (sel) -- (l0);
    \draw[seta] (sel.south) to[out=250, in=290]
        node[rot, below=2pt] {atualiza os termos cobertos; repete até $n_i$ instâncias}
        (cand.south);

    % ---------------- avaliação (sem afirmar resultado) --------------------
    \node[rot, anchor=west, text width=124mm, align=left] (tit3) at (-1.9,-8.6)
        {\textbf{Avaliação} (resultados na Seção~\ref{sec:res-l0-drisl}): o
         $L_0$ do DRI-SL é comparado ao da amostragem aleatória e ao envelope
         de desempenho do algoritmo genético, que atua como limite empírico
         de referência.};
\end{tikzpicture}

%   A citação do codificador de sentenças (\citep{Reimers2019SBERT}) vai na
%   LEGENDA da figure, não pintada no desenho — assim o check-bib a enxerga.
% Uso standalone (pré-visualização):
%   \documentclass[tikz,border=6pt]{standalone}
%   \usetikzlibrary{arrows.meta, positioning}
%   \begin{document}% ============================================================================
% ESQUEMA PROPOSTO (banca, sugestão — não mergeado): o algoritmo DRI-SL como
% fluxo em duas fases — (i) densidade semântica e (ii) variedade lexical
% intragrupo — seguido da avaliação. Ilustra a Seção sec:metodo-drisl.
%
% Uso no Cap. 3 (dentro de um ambiente figure):
%   % ============================================================================
% ESQUEMA PROPOSTO (banca, sugestão — não mergeado): o algoritmo DRI-SL como
% fluxo em duas fases — (i) densidade semântica e (ii) variedade lexical
% intragrupo — seguido da avaliação. Ilustra a Seção sec:metodo-drisl.
%
% Uso no Cap. 3 (dentro de um ambiente figure):
%   % ============================================================================
% ESQUEMA PROPOSTO (banca, sugestão — não mergeado): o algoritmo DRI-SL como
% fluxo em duas fases — (i) densidade semântica e (ii) variedade lexical
% intragrupo — seguido da avaliação. Ilustra a Seção sec:metodo-drisl.
%
% Uso no Cap. 3 (dentro de um ambiente figure):
%   \input{3-metodo/esquemas-propostos/esq-drisl}
%   A citação do codificador de sentenças (\citep{Reimers2019SBERT}) vai na
%   LEGENDA da figure, não pintada no desenho — assim o check-bib a enxerga.
% Uso standalone (pré-visualização):
%   \documentclass[tikz,border=6pt]{standalone}
%   \usetikzlibrary{arrows.meta, positioning}
%   \begin{document}\input{esq-drisl.tex}\end{document}
%   (no modo standalone, troque as \ref{...} por texto fixo ou defina-as)
% Requer apenas arrows.meta e positioning (já em packages.tex). Legível em
% P&B: etapa de medição/entrada = cinza claro; processamento = branco;
% produto = cinza médio; remissão = texto simples.
% FREEZE respeitado: nenhum número de resultado; só os símbolos da própria
% seção (U0, L0, I, Nc, Gi, ni) e nenhum gráfico decorativo que pareça dado.
% Terminologia em camadas: a figura não afirma o resultado da comparação
% (isso é do Capítulo 4); só diz CONTRA O QUE o DRI-SL é avaliado.
% ============================================================================
\begin{tikzpicture}[
    x=1cm, y=1cm,
    caixa/.style={draw, rounded corners=2pt, align=center, text width=34mm,
                  minimum height=11mm, inner sep=3pt, font=\footnotesize},
    entrada/.style={caixa, fill=gray!8},
    produto/.style={caixa, fill=gray!14},
    seta/.style={-{Stealth[length=2.4mm]}, thick},
    rot/.style={font=\scriptsize, align=center, inner sep=2pt}
]
    % ---------------- fase (i): densidade semântica ------------------------
    \node[rot, anchor=west] (tit1) at (-1.9,1.1)
        {\textbf{Fase (i) --- densidade semântica}};

    \node[entrada] (u0) at (0,0)
        {conjunto não rotulado $U_0$ (textos curtos)};
    \node[caixa, text width=36mm] (enc) at (4.6,0)
        {codificador de\\sentenças: projeção\\no espaço semântico};
    \node[caixa] (km) at (9.2,0)
        {$k$-médias em $N_c$ grupos semânticos};
    \draw[seta] (u0) -- (enc);
    \draw[seta] (enc) -- (km);

    \node[produto] (aloc) at (9.2,-2.3)
        {cotas proporcionais ao tamanho de cada grupo:
         $\sum_{i=1}^{N_c} n_i = I$};
    \draw[seta] (km) -- (aloc);

    % ---------------- fase (ii): variedade lexical intragrupo --------------
    % o título da fase É o rótulo da seta em L (um elemento só: título de
    % banda separado era atravessado pela descida da seta — achado visual)
    \node[caixa] (cand) at (0,-5.4)
        {candidatos do grupo $G_i$ (ainda não selecionados)};
    \node[caixa, text width=40mm] (sel) at (4.6,-5.4)
        {seleciona a instância de maior escore de novidade, ponderado pelo
         perfil TF--IDF do grupo};
    \node[produto, text width=44mm] (l0) at (9.6,-5.4)
        {conjunto inicial $L_0$ (\mbox{$|L_0|=I$}):\\
         diversidade semântica entre\\
         grupos $+$ variedade lexical\\
         dentro de cada grupo};

    % a linha em L conduz da alocação para o laço de seleção; o rótulo na
    % horizontal é o título da fase (desenhada aqui, depois dos nós)
    \draw[seta] (aloc.south) -- ++(0,-0.55) -| (cand.north);
    % rótulo-título posicionado à mão sobre o trecho horizontal, com as
    % bordas afastadas das duas descidas verticais (x=0 e x=9,2)
    \node[rot, anchor=south] at (4.2,-3.42)
        {\textbf{Fase (ii) --- variedade lexical intragrupo}:\\
         para cada grupo $G_i$, com cota $n_i$};

    \draw[seta] (cand) -- (sel);
    \draw[seta] (sel) -- (l0);
    \draw[seta] (sel.south) to[out=250, in=290]
        node[rot, below=2pt] {atualiza os termos cobertos; repete até $n_i$ instâncias}
        (cand.south);

    % ---------------- avaliação (sem afirmar resultado) --------------------
    \node[rot, anchor=west, text width=124mm, align=left] (tit3) at (-1.9,-8.6)
        {\textbf{Avaliação} (resultados na Seção~\ref{sec:res-l0-drisl}): o
         $L_0$ do DRI-SL é comparado ao da amostragem aleatória e ao envelope
         de desempenho do algoritmo genético, que atua como limite empírico
         de referência.};
\end{tikzpicture}

%   A citação do codificador de sentenças (\citep{Reimers2019SBERT}) vai na
%   LEGENDA da figure, não pintada no desenho — assim o check-bib a enxerga.
% Uso standalone (pré-visualização):
%   \documentclass[tikz,border=6pt]{standalone}
%   \usetikzlibrary{arrows.meta, positioning}
%   \begin{document}% ============================================================================
% ESQUEMA PROPOSTO (banca, sugestão — não mergeado): o algoritmo DRI-SL como
% fluxo em duas fases — (i) densidade semântica e (ii) variedade lexical
% intragrupo — seguido da avaliação. Ilustra a Seção sec:metodo-drisl.
%
% Uso no Cap. 3 (dentro de um ambiente figure):
%   \input{3-metodo/esquemas-propostos/esq-drisl}
%   A citação do codificador de sentenças (\citep{Reimers2019SBERT}) vai na
%   LEGENDA da figure, não pintada no desenho — assim o check-bib a enxerga.
% Uso standalone (pré-visualização):
%   \documentclass[tikz,border=6pt]{standalone}
%   \usetikzlibrary{arrows.meta, positioning}
%   \begin{document}\input{esq-drisl.tex}\end{document}
%   (no modo standalone, troque as \ref{...} por texto fixo ou defina-as)
% Requer apenas arrows.meta e positioning (já em packages.tex). Legível em
% P&B: etapa de medição/entrada = cinza claro; processamento = branco;
% produto = cinza médio; remissão = texto simples.
% FREEZE respeitado: nenhum número de resultado; só os símbolos da própria
% seção (U0, L0, I, Nc, Gi, ni) e nenhum gráfico decorativo que pareça dado.
% Terminologia em camadas: a figura não afirma o resultado da comparação
% (isso é do Capítulo 4); só diz CONTRA O QUE o DRI-SL é avaliado.
% ============================================================================
\begin{tikzpicture}[
    x=1cm, y=1cm,
    caixa/.style={draw, rounded corners=2pt, align=center, text width=34mm,
                  minimum height=11mm, inner sep=3pt, font=\footnotesize},
    entrada/.style={caixa, fill=gray!8},
    produto/.style={caixa, fill=gray!14},
    seta/.style={-{Stealth[length=2.4mm]}, thick},
    rot/.style={font=\scriptsize, align=center, inner sep=2pt}
]
    % ---------------- fase (i): densidade semântica ------------------------
    \node[rot, anchor=west] (tit1) at (-1.9,1.1)
        {\textbf{Fase (i) --- densidade semântica}};

    \node[entrada] (u0) at (0,0)
        {conjunto não rotulado $U_0$ (textos curtos)};
    \node[caixa, text width=36mm] (enc) at (4.6,0)
        {codificador de\\sentenças: projeção\\no espaço semântico};
    \node[caixa] (km) at (9.2,0)
        {$k$-médias em $N_c$ grupos semânticos};
    \draw[seta] (u0) -- (enc);
    \draw[seta] (enc) -- (km);

    \node[produto] (aloc) at (9.2,-2.3)
        {cotas proporcionais ao tamanho de cada grupo:
         $\sum_{i=1}^{N_c} n_i = I$};
    \draw[seta] (km) -- (aloc);

    % ---------------- fase (ii): variedade lexical intragrupo --------------
    % o título da fase É o rótulo da seta em L (um elemento só: título de
    % banda separado era atravessado pela descida da seta — achado visual)
    \node[caixa] (cand) at (0,-5.4)
        {candidatos do grupo $G_i$ (ainda não selecionados)};
    \node[caixa, text width=40mm] (sel) at (4.6,-5.4)
        {seleciona a instância de maior escore de novidade, ponderado pelo
         perfil TF--IDF do grupo};
    \node[produto, text width=44mm] (l0) at (9.6,-5.4)
        {conjunto inicial $L_0$ (\mbox{$|L_0|=I$}):\\
         diversidade semântica entre\\
         grupos $+$ variedade lexical\\
         dentro de cada grupo};

    % a linha em L conduz da alocação para o laço de seleção; o rótulo na
    % horizontal é o título da fase (desenhada aqui, depois dos nós)
    \draw[seta] (aloc.south) -- ++(0,-0.55) -| (cand.north);
    % rótulo-título posicionado à mão sobre o trecho horizontal, com as
    % bordas afastadas das duas descidas verticais (x=0 e x=9,2)
    \node[rot, anchor=south] at (4.2,-3.42)
        {\textbf{Fase (ii) --- variedade lexical intragrupo}:\\
         para cada grupo $G_i$, com cota $n_i$};

    \draw[seta] (cand) -- (sel);
    \draw[seta] (sel) -- (l0);
    \draw[seta] (sel.south) to[out=250, in=290]
        node[rot, below=2pt] {atualiza os termos cobertos; repete até $n_i$ instâncias}
        (cand.south);

    % ---------------- avaliação (sem afirmar resultado) --------------------
    \node[rot, anchor=west, text width=124mm, align=left] (tit3) at (-1.9,-8.6)
        {\textbf{Avaliação} (resultados na Seção~\ref{sec:res-l0-drisl}): o
         $L_0$ do DRI-SL é comparado ao da amostragem aleatória e ao envelope
         de desempenho do algoritmo genético, que atua como limite empírico
         de referência.};
\end{tikzpicture}
\end{document}
%   (no modo standalone, troque as \ref{...} por texto fixo ou defina-as)
% Requer apenas arrows.meta e positioning (já em packages.tex). Legível em
% P&B: etapa de medição/entrada = cinza claro; processamento = branco;
% produto = cinza médio; remissão = texto simples.
% FREEZE respeitado: nenhum número de resultado; só os símbolos da própria
% seção (U0, L0, I, Nc, Gi, ni) e nenhum gráfico decorativo que pareça dado.
% Terminologia em camadas: a figura não afirma o resultado da comparação
% (isso é do Capítulo 4); só diz CONTRA O QUE o DRI-SL é avaliado.
% ============================================================================
\begin{tikzpicture}[
    x=1cm, y=1cm,
    caixa/.style={draw, rounded corners=2pt, align=center, text width=34mm,
                  minimum height=11mm, inner sep=3pt, font=\footnotesize},
    entrada/.style={caixa, fill=gray!8},
    produto/.style={caixa, fill=gray!14},
    seta/.style={-{Stealth[length=2.4mm]}, thick},
    rot/.style={font=\scriptsize, align=center, inner sep=2pt}
]
    % ---------------- fase (i): densidade semântica ------------------------
    \node[rot, anchor=west] (tit1) at (-1.9,1.1)
        {\textbf{Fase (i) --- densidade semântica}};

    \node[entrada] (u0) at (0,0)
        {conjunto não rotulado $U_0$ (textos curtos)};
    \node[caixa, text width=36mm] (enc) at (4.6,0)
        {codificador de\\sentenças: projeção\\no espaço semântico};
    \node[caixa] (km) at (9.2,0)
        {$k$-médias em $N_c$ grupos semânticos};
    \draw[seta] (u0) -- (enc);
    \draw[seta] (enc) -- (km);

    \node[produto] (aloc) at (9.2,-2.3)
        {cotas proporcionais ao tamanho de cada grupo:
         $\sum_{i=1}^{N_c} n_i = I$};
    \draw[seta] (km) -- (aloc);

    % ---------------- fase (ii): variedade lexical intragrupo --------------
    % o título da fase É o rótulo da seta em L (um elemento só: título de
    % banda separado era atravessado pela descida da seta — achado visual)
    \node[caixa] (cand) at (0,-5.4)
        {candidatos do grupo $G_i$ (ainda não selecionados)};
    \node[caixa, text width=40mm] (sel) at (4.6,-5.4)
        {seleciona a instância de maior escore de novidade, ponderado pelo
         perfil TF--IDF do grupo};
    \node[produto, text width=44mm] (l0) at (9.6,-5.4)
        {conjunto inicial $L_0$ (\mbox{$|L_0|=I$}):\\
         diversidade semântica entre\\
         grupos $+$ variedade lexical\\
         dentro de cada grupo};

    % a linha em L conduz da alocação para o laço de seleção; o rótulo na
    % horizontal é o título da fase (desenhada aqui, depois dos nós)
    \draw[seta] (aloc.south) -- ++(0,-0.55) -| (cand.north);
    % rótulo-título posicionado à mão sobre o trecho horizontal, com as
    % bordas afastadas das duas descidas verticais (x=0 e x=9,2)
    \node[rot, anchor=south] at (4.2,-3.42)
        {\textbf{Fase (ii) --- variedade lexical intragrupo}:\\
         para cada grupo $G_i$, com cota $n_i$};

    \draw[seta] (cand) -- (sel);
    \draw[seta] (sel) -- (l0);
    \draw[seta] (sel.south) to[out=250, in=290]
        node[rot, below=2pt] {atualiza os termos cobertos; repete até $n_i$ instâncias}
        (cand.south);

    % ---------------- avaliação (sem afirmar resultado) --------------------
    \node[rot, anchor=west, text width=124mm, align=left] (tit3) at (-1.9,-8.6)
        {\textbf{Avaliação} (resultados na Seção~\ref{sec:res-l0-drisl}): o
         $L_0$ do DRI-SL é comparado ao da amostragem aleatória e ao envelope
         de desempenho do algoritmo genético, que atua como limite empírico
         de referência.};
\end{tikzpicture}

%   A citação do codificador de sentenças (\citep{Reimers2019SBERT}) vai na
%   LEGENDA da figure, não pintada no desenho — assim o check-bib a enxerga.
% Uso standalone (pré-visualização):
%   \documentclass[tikz,border=6pt]{standalone}
%   \usetikzlibrary{arrows.meta, positioning}
%   \begin{document}% ============================================================================
% ESQUEMA PROPOSTO (banca, sugestão — não mergeado): o algoritmo DRI-SL como
% fluxo em duas fases — (i) densidade semântica e (ii) variedade lexical
% intragrupo — seguido da avaliação. Ilustra a Seção sec:metodo-drisl.
%
% Uso no Cap. 3 (dentro de um ambiente figure):
%   % ============================================================================
% ESQUEMA PROPOSTO (banca, sugestão — não mergeado): o algoritmo DRI-SL como
% fluxo em duas fases — (i) densidade semântica e (ii) variedade lexical
% intragrupo — seguido da avaliação. Ilustra a Seção sec:metodo-drisl.
%
% Uso no Cap. 3 (dentro de um ambiente figure):
%   \input{3-metodo/esquemas-propostos/esq-drisl}
%   A citação do codificador de sentenças (\citep{Reimers2019SBERT}) vai na
%   LEGENDA da figure, não pintada no desenho — assim o check-bib a enxerga.
% Uso standalone (pré-visualização):
%   \documentclass[tikz,border=6pt]{standalone}
%   \usetikzlibrary{arrows.meta, positioning}
%   \begin{document}\input{esq-drisl.tex}\end{document}
%   (no modo standalone, troque as \ref{...} por texto fixo ou defina-as)
% Requer apenas arrows.meta e positioning (já em packages.tex). Legível em
% P&B: etapa de medição/entrada = cinza claro; processamento = branco;
% produto = cinza médio; remissão = texto simples.
% FREEZE respeitado: nenhum número de resultado; só os símbolos da própria
% seção (U0, L0, I, Nc, Gi, ni) e nenhum gráfico decorativo que pareça dado.
% Terminologia em camadas: a figura não afirma o resultado da comparação
% (isso é do Capítulo 4); só diz CONTRA O QUE o DRI-SL é avaliado.
% ============================================================================
\begin{tikzpicture}[
    x=1cm, y=1cm,
    caixa/.style={draw, rounded corners=2pt, align=center, text width=34mm,
                  minimum height=11mm, inner sep=3pt, font=\footnotesize},
    entrada/.style={caixa, fill=gray!8},
    produto/.style={caixa, fill=gray!14},
    seta/.style={-{Stealth[length=2.4mm]}, thick},
    rot/.style={font=\scriptsize, align=center, inner sep=2pt}
]
    % ---------------- fase (i): densidade semântica ------------------------
    \node[rot, anchor=west] (tit1) at (-1.9,1.1)
        {\textbf{Fase (i) --- densidade semântica}};

    \node[entrada] (u0) at (0,0)
        {conjunto não rotulado $U_0$ (textos curtos)};
    \node[caixa, text width=36mm] (enc) at (4.6,0)
        {codificador de\\sentenças: projeção\\no espaço semântico};
    \node[caixa] (km) at (9.2,0)
        {$k$-médias em $N_c$ grupos semânticos};
    \draw[seta] (u0) -- (enc);
    \draw[seta] (enc) -- (km);

    \node[produto] (aloc) at (9.2,-2.3)
        {cotas proporcionais ao tamanho de cada grupo:
         $\sum_{i=1}^{N_c} n_i = I$};
    \draw[seta] (km) -- (aloc);

    % ---------------- fase (ii): variedade lexical intragrupo --------------
    % o título da fase É o rótulo da seta em L (um elemento só: título de
    % banda separado era atravessado pela descida da seta — achado visual)
    \node[caixa] (cand) at (0,-5.4)
        {candidatos do grupo $G_i$ (ainda não selecionados)};
    \node[caixa, text width=40mm] (sel) at (4.6,-5.4)
        {seleciona a instância de maior escore de novidade, ponderado pelo
         perfil TF--IDF do grupo};
    \node[produto, text width=44mm] (l0) at (9.6,-5.4)
        {conjunto inicial $L_0$ (\mbox{$|L_0|=I$}):\\
         diversidade semântica entre\\
         grupos $+$ variedade lexical\\
         dentro de cada grupo};

    % a linha em L conduz da alocação para o laço de seleção; o rótulo na
    % horizontal é o título da fase (desenhada aqui, depois dos nós)
    \draw[seta] (aloc.south) -- ++(0,-0.55) -| (cand.north);
    % rótulo-título posicionado à mão sobre o trecho horizontal, com as
    % bordas afastadas das duas descidas verticais (x=0 e x=9,2)
    \node[rot, anchor=south] at (4.2,-3.42)
        {\textbf{Fase (ii) --- variedade lexical intragrupo}:\\
         para cada grupo $G_i$, com cota $n_i$};

    \draw[seta] (cand) -- (sel);
    \draw[seta] (sel) -- (l0);
    \draw[seta] (sel.south) to[out=250, in=290]
        node[rot, below=2pt] {atualiza os termos cobertos; repete até $n_i$ instâncias}
        (cand.south);

    % ---------------- avaliação (sem afirmar resultado) --------------------
    \node[rot, anchor=west, text width=124mm, align=left] (tit3) at (-1.9,-8.6)
        {\textbf{Avaliação} (resultados na Seção~\ref{sec:res-l0-drisl}): o
         $L_0$ do DRI-SL é comparado ao da amostragem aleatória e ao envelope
         de desempenho do algoritmo genético, que atua como limite empírico
         de referência.};
\end{tikzpicture}

%   A citação do codificador de sentenças (\citep{Reimers2019SBERT}) vai na
%   LEGENDA da figure, não pintada no desenho — assim o check-bib a enxerga.
% Uso standalone (pré-visualização):
%   \documentclass[tikz,border=6pt]{standalone}
%   \usetikzlibrary{arrows.meta, positioning}
%   \begin{document}% ============================================================================
% ESQUEMA PROPOSTO (banca, sugestão — não mergeado): o algoritmo DRI-SL como
% fluxo em duas fases — (i) densidade semântica e (ii) variedade lexical
% intragrupo — seguido da avaliação. Ilustra a Seção sec:metodo-drisl.
%
% Uso no Cap. 3 (dentro de um ambiente figure):
%   \input{3-metodo/esquemas-propostos/esq-drisl}
%   A citação do codificador de sentenças (\citep{Reimers2019SBERT}) vai na
%   LEGENDA da figure, não pintada no desenho — assim o check-bib a enxerga.
% Uso standalone (pré-visualização):
%   \documentclass[tikz,border=6pt]{standalone}
%   \usetikzlibrary{arrows.meta, positioning}
%   \begin{document}\input{esq-drisl.tex}\end{document}
%   (no modo standalone, troque as \ref{...} por texto fixo ou defina-as)
% Requer apenas arrows.meta e positioning (já em packages.tex). Legível em
% P&B: etapa de medição/entrada = cinza claro; processamento = branco;
% produto = cinza médio; remissão = texto simples.
% FREEZE respeitado: nenhum número de resultado; só os símbolos da própria
% seção (U0, L0, I, Nc, Gi, ni) e nenhum gráfico decorativo que pareça dado.
% Terminologia em camadas: a figura não afirma o resultado da comparação
% (isso é do Capítulo 4); só diz CONTRA O QUE o DRI-SL é avaliado.
% ============================================================================
\begin{tikzpicture}[
    x=1cm, y=1cm,
    caixa/.style={draw, rounded corners=2pt, align=center, text width=34mm,
                  minimum height=11mm, inner sep=3pt, font=\footnotesize},
    entrada/.style={caixa, fill=gray!8},
    produto/.style={caixa, fill=gray!14},
    seta/.style={-{Stealth[length=2.4mm]}, thick},
    rot/.style={font=\scriptsize, align=center, inner sep=2pt}
]
    % ---------------- fase (i): densidade semântica ------------------------
    \node[rot, anchor=west] (tit1) at (-1.9,1.1)
        {\textbf{Fase (i) --- densidade semântica}};

    \node[entrada] (u0) at (0,0)
        {conjunto não rotulado $U_0$ (textos curtos)};
    \node[caixa, text width=36mm] (enc) at (4.6,0)
        {codificador de\\sentenças: projeção\\no espaço semântico};
    \node[caixa] (km) at (9.2,0)
        {$k$-médias em $N_c$ grupos semânticos};
    \draw[seta] (u0) -- (enc);
    \draw[seta] (enc) -- (km);

    \node[produto] (aloc) at (9.2,-2.3)
        {cotas proporcionais ao tamanho de cada grupo:
         $\sum_{i=1}^{N_c} n_i = I$};
    \draw[seta] (km) -- (aloc);

    % ---------------- fase (ii): variedade lexical intragrupo --------------
    % o título da fase É o rótulo da seta em L (um elemento só: título de
    % banda separado era atravessado pela descida da seta — achado visual)
    \node[caixa] (cand) at (0,-5.4)
        {candidatos do grupo $G_i$ (ainda não selecionados)};
    \node[caixa, text width=40mm] (sel) at (4.6,-5.4)
        {seleciona a instância de maior escore de novidade, ponderado pelo
         perfil TF--IDF do grupo};
    \node[produto, text width=44mm] (l0) at (9.6,-5.4)
        {conjunto inicial $L_0$ (\mbox{$|L_0|=I$}):\\
         diversidade semântica entre\\
         grupos $+$ variedade lexical\\
         dentro de cada grupo};

    % a linha em L conduz da alocação para o laço de seleção; o rótulo na
    % horizontal é o título da fase (desenhada aqui, depois dos nós)
    \draw[seta] (aloc.south) -- ++(0,-0.55) -| (cand.north);
    % rótulo-título posicionado à mão sobre o trecho horizontal, com as
    % bordas afastadas das duas descidas verticais (x=0 e x=9,2)
    \node[rot, anchor=south] at (4.2,-3.42)
        {\textbf{Fase (ii) --- variedade lexical intragrupo}:\\
         para cada grupo $G_i$, com cota $n_i$};

    \draw[seta] (cand) -- (sel);
    \draw[seta] (sel) -- (l0);
    \draw[seta] (sel.south) to[out=250, in=290]
        node[rot, below=2pt] {atualiza os termos cobertos; repete até $n_i$ instâncias}
        (cand.south);

    % ---------------- avaliação (sem afirmar resultado) --------------------
    \node[rot, anchor=west, text width=124mm, align=left] (tit3) at (-1.9,-8.6)
        {\textbf{Avaliação} (resultados na Seção~\ref{sec:res-l0-drisl}): o
         $L_0$ do DRI-SL é comparado ao da amostragem aleatória e ao envelope
         de desempenho do algoritmo genético, que atua como limite empírico
         de referência.};
\end{tikzpicture}
\end{document}
%   (no modo standalone, troque as \ref{...} por texto fixo ou defina-as)
% Requer apenas arrows.meta e positioning (já em packages.tex). Legível em
% P&B: etapa de medição/entrada = cinza claro; processamento = branco;
% produto = cinza médio; remissão = texto simples.
% FREEZE respeitado: nenhum número de resultado; só os símbolos da própria
% seção (U0, L0, I, Nc, Gi, ni) e nenhum gráfico decorativo que pareça dado.
% Terminologia em camadas: a figura não afirma o resultado da comparação
% (isso é do Capítulo 4); só diz CONTRA O QUE o DRI-SL é avaliado.
% ============================================================================
\begin{tikzpicture}[
    x=1cm, y=1cm,
    caixa/.style={draw, rounded corners=2pt, align=center, text width=34mm,
                  minimum height=11mm, inner sep=3pt, font=\footnotesize},
    entrada/.style={caixa, fill=gray!8},
    produto/.style={caixa, fill=gray!14},
    seta/.style={-{Stealth[length=2.4mm]}, thick},
    rot/.style={font=\scriptsize, align=center, inner sep=2pt}
]
    % ---------------- fase (i): densidade semântica ------------------------
    \node[rot, anchor=west] (tit1) at (-1.9,1.1)
        {\textbf{Fase (i) --- densidade semântica}};

    \node[entrada] (u0) at (0,0)
        {conjunto não rotulado $U_0$ (textos curtos)};
    \node[caixa, text width=36mm] (enc) at (4.6,0)
        {codificador de\\sentenças: projeção\\no espaço semântico};
    \node[caixa] (km) at (9.2,0)
        {$k$-médias em $N_c$ grupos semânticos};
    \draw[seta] (u0) -- (enc);
    \draw[seta] (enc) -- (km);

    \node[produto] (aloc) at (9.2,-2.3)
        {cotas proporcionais ao tamanho de cada grupo:
         $\sum_{i=1}^{N_c} n_i = I$};
    \draw[seta] (km) -- (aloc);

    % ---------------- fase (ii): variedade lexical intragrupo --------------
    % o título da fase É o rótulo da seta em L (um elemento só: título de
    % banda separado era atravessado pela descida da seta — achado visual)
    \node[caixa] (cand) at (0,-5.4)
        {candidatos do grupo $G_i$ (ainda não selecionados)};
    \node[caixa, text width=40mm] (sel) at (4.6,-5.4)
        {seleciona a instância de maior escore de novidade, ponderado pelo
         perfil TF--IDF do grupo};
    \node[produto, text width=44mm] (l0) at (9.6,-5.4)
        {conjunto inicial $L_0$ (\mbox{$|L_0|=I$}):\\
         diversidade semântica entre\\
         grupos $+$ variedade lexical\\
         dentro de cada grupo};

    % a linha em L conduz da alocação para o laço de seleção; o rótulo na
    % horizontal é o título da fase (desenhada aqui, depois dos nós)
    \draw[seta] (aloc.south) -- ++(0,-0.55) -| (cand.north);
    % rótulo-título posicionado à mão sobre o trecho horizontal, com as
    % bordas afastadas das duas descidas verticais (x=0 e x=9,2)
    \node[rot, anchor=south] at (4.2,-3.42)
        {\textbf{Fase (ii) --- variedade lexical intragrupo}:\\
         para cada grupo $G_i$, com cota $n_i$};

    \draw[seta] (cand) -- (sel);
    \draw[seta] (sel) -- (l0);
    \draw[seta] (sel.south) to[out=250, in=290]
        node[rot, below=2pt] {atualiza os termos cobertos; repete até $n_i$ instâncias}
        (cand.south);

    % ---------------- avaliação (sem afirmar resultado) --------------------
    \node[rot, anchor=west, text width=124mm, align=left] (tit3) at (-1.9,-8.6)
        {\textbf{Avaliação} (resultados na Seção~\ref{sec:res-l0-drisl}): o
         $L_0$ do DRI-SL é comparado ao da amostragem aleatória e ao envelope
         de desempenho do algoritmo genético, que atua como limite empírico
         de referência.};
\end{tikzpicture}
\end{document}
%   (no modo standalone, troque as \ref{...} por texto fixo ou defina-as)
% Requer apenas arrows.meta e positioning (já em packages.tex). Legível em
% P&B: etapa de medição/entrada = cinza claro; processamento = branco;
% produto = cinza médio; remissão = texto simples.
% FREEZE respeitado: nenhum número de resultado; só os símbolos da própria
% seção (U0, L0, I, Nc, Gi, ni) e nenhum gráfico decorativo que pareça dado.
% Terminologia em camadas: a figura não afirma o resultado da comparação
% (isso é do Capítulo 4); só diz CONTRA O QUE o DRI-SL é avaliado.
% ============================================================================
\begin{tikzpicture}[
    x=1cm, y=1cm,
    caixa/.style={draw, rounded corners=2pt, align=center, text width=34mm,
                  minimum height=11mm, inner sep=3pt, font=\footnotesize},
    entrada/.style={caixa, fill=gray!8},
    produto/.style={caixa, fill=gray!14},
    seta/.style={-{Stealth[length=2.4mm]}, thick},
    rot/.style={font=\scriptsize, align=center, inner sep=2pt}
]
    % ---------------- fase (i): densidade semântica ------------------------
    \node[rot, anchor=west] (tit1) at (-1.9,1.1)
        {\textbf{Fase (i) --- densidade semântica}};

    \node[entrada] (u0) at (0,0)
        {conjunto não rotulado $U_0$ (textos curtos)};
    \node[caixa, text width=36mm] (enc) at (4.6,0)
        {codificador de\\sentenças: projeção\\no espaço semântico};
    \node[caixa] (km) at (9.2,0)
        {$k$-médias em $N_c$ grupos semânticos};
    \draw[seta] (u0) -- (enc);
    \draw[seta] (enc) -- (km);

    \node[produto] (aloc) at (9.2,-2.3)
        {cotas proporcionais ao tamanho de cada grupo:
         $\sum_{i=1}^{N_c} n_i = I$};
    \draw[seta] (km) -- (aloc);

    % ---------------- fase (ii): variedade lexical intragrupo --------------
    % o título da fase É o rótulo da seta em L (um elemento só: título de
    % banda separado era atravessado pela descida da seta — achado visual)
    \node[caixa] (cand) at (0,-5.4)
        {candidatos do grupo $G_i$ (ainda não selecionados)};
    \node[caixa, text width=40mm] (sel) at (4.6,-5.4)
        {seleciona a instância de maior escore de novidade, ponderado pelo
         perfil TF--IDF do grupo};
    \node[produto, text width=44mm] (l0) at (9.6,-5.4)
        {conjunto inicial $L_0$ (\mbox{$|L_0|=I$}):\\
         diversidade semântica entre\\
         grupos $+$ variedade lexical\\
         dentro de cada grupo};

    % a linha em L conduz da alocação para o laço de seleção; o rótulo na
    % horizontal é o título da fase (desenhada aqui, depois dos nós)
    \draw[seta] (aloc.south) -- ++(0,-0.55) -| (cand.north);
    % rótulo-título posicionado à mão sobre o trecho horizontal, com as
    % bordas afastadas das duas descidas verticais (x=0 e x=9,2)
    \node[rot, anchor=south] at (4.2,-3.42)
        {\textbf{Fase (ii) --- variedade lexical intragrupo}:\\
         para cada grupo $G_i$, com cota $n_i$};

    \draw[seta] (cand) -- (sel);
    \draw[seta] (sel) -- (l0);
    \draw[seta] (sel.south) to[out=250, in=290]
        node[rot, below=2pt] {atualiza os termos cobertos; repete até $n_i$ instâncias}
        (cand.south);

    % ---------------- avaliação (sem afirmar resultado) --------------------
    \node[rot, anchor=west, text width=124mm, align=left] (tit3) at (-1.9,-8.6)
        {\textbf{Avaliação} (resultados na Seção~\ref{sec:res-l0-drisl}): o
         $L_0$ do DRI-SL é comparado ao da amostragem aleatória e ao envelope
         de desempenho do algoritmo genético, que atua como limite empírico
         de referência.};
\end{tikzpicture}
\end{document}
%   (no modo standalone, troque as \ref{...} por texto fixo ou defina-as)
% Requer apenas arrows.meta e positioning (já em packages.tex). Legível em
% P&B: etapa de medição/entrada = cinza claro; processamento = branco;
% produto = cinza médio; remissão = texto simples.
% FREEZE respeitado: nenhum número de resultado; só os símbolos da própria
% seção (U0, L0, I, Nc, Gi, ni) e nenhum gráfico decorativo que pareça dado.
% Terminologia em camadas: a figura não afirma o resultado da comparação
% (isso é do Capítulo 4); só diz CONTRA O QUE o DRI-SL é avaliado.
% ============================================================================
\begin{tikzpicture}[
    x=1cm, y=1cm,
    caixa/.style={draw, rounded corners=2pt, align=center, text width=34mm,
                  minimum height=11mm, inner sep=3pt, font=\footnotesize},
    entrada/.style={caixa, fill=gray!8},
    produto/.style={caixa, fill=gray!14},
    seta/.style={-{Stealth[length=2.4mm]}, thick},
    rot/.style={font=\scriptsize, align=center, inner sep=2pt}
]
    % ---------------- fase (i): densidade semântica ------------------------
    \node[rot, anchor=west] (tit1) at (-1.9,1.1)
        {\textbf{Fase (i) --- densidade semântica}};

    \node[entrada] (u0) at (0,0)
        {conjunto não rotulado $U_0$ (textos curtos)};
    \node[caixa, text width=36mm] (enc) at (4.6,0)
        {codificador de\\sentenças: projeção\\no espaço semântico};
    \node[caixa] (km) at (9.2,0)
        {$k$-médias em $N_c$ grupos semânticos};
    \draw[seta] (u0) -- (enc);
    \draw[seta] (enc) -- (km);

    \node[produto] (aloc) at (9.2,-2.3)
        {cotas proporcionais ao tamanho de cada grupo:
         $\sum_{i=1}^{N_c} n_i = I$};
    \draw[seta] (km) -- (aloc);

    % ---------------- fase (ii): variedade lexical intragrupo --------------
    % o título da fase É o rótulo da seta em L (um elemento só: título de
    % banda separado era atravessado pela descida da seta — achado visual)
    \node[caixa] (cand) at (0,-5.4)
        {candidatos do grupo $G_i$ (ainda não selecionados)};
    \node[caixa, text width=40mm] (sel) at (4.6,-5.4)
        {seleciona a instância de maior escore de novidade, ponderado pelo
         perfil TF--IDF do grupo};
    \node[produto, text width=44mm] (l0) at (9.6,-5.4)
        {conjunto inicial $L_0$ (\mbox{$|L_0|=I$}):\\
         diversidade semântica entre\\
         grupos $+$ variedade lexical\\
         dentro de cada grupo};

    % a linha em L conduz da alocação para o laço de seleção; o rótulo na
    % horizontal é o título da fase (desenhada aqui, depois dos nós)
    \draw[seta] (aloc.south) -- ++(0,-0.55) -| (cand.north);
    % rótulo-título posicionado à mão sobre o trecho horizontal, com as
    % bordas afastadas das duas descidas verticais (x=0 e x=9,2)
    \node[rot, anchor=south] at (4.2,-3.42)
        {\textbf{Fase (ii) --- variedade lexical intragrupo}:\\
         para cada grupo $G_i$, com cota $n_i$};

    \draw[seta] (cand) -- (sel);
    \draw[seta] (sel) -- (l0);
    \draw[seta] (sel.south) to[out=250, in=290]
        node[rot, below=2pt] {atualiza os termos cobertos; repete até $n_i$ instâncias}
        (cand.south);

    % ---------------- avaliação (sem afirmar resultado) --------------------
    \node[rot, anchor=west, text width=124mm, align=left] (tit3) at (-1.9,-8.6)
        {\textbf{Avaliação} (resultados na Seção~\ref{sec:res-l0-drisl}): o
         $L_0$ do DRI-SL é comparado ao da amostragem aleatória e ao envelope
         de desempenho do algoritmo genético, que atua como limite empírico
         de referência.};
\end{tikzpicture}
